\documentclass[12pt]{article}

\usepackage[a4paper, top=2cm, bottom=2cm, left=2.5cm, right=2.5cm]{geometry}

\usepackage{amsmath, amssymb}
\usepackage{tikz}
\usetikzlibrary{arrows.meta,positioning,fit,calc}
\usepackage{multirow}
\usepackage{array}
\usepackage{siunitx}
\usepackage{enumitem}
\usepackage{float}
\usepackage{fontspec}
\usetikzlibrary{decorations.pathmorphing}
\usepackage{pdflscape}
\usepackage{tabularx}
\usepackage{booktabs}
\usepackage{graphicx}
\usepackage{listings}
\usepackage[table]{xcolor}
\usepackage{hyperref}
\usepackage{bookmark}
\usepackage{pdfpages}

% ---------------- FONTS ----------------
\setmainfont[
    UprightFont = Handlee-Regular.ttf,
    AutoFakeBold = 2.5
]{Handlee}

\newfontfamily\headingfont{PT.ttf}

\setmonofont{DejaVu Sans Mono}

% ---------------- HANDWRITTEN MATH ----------------
\usepackage{mathastext}
\Mathastext
\everymath{\displaystyle}

% ---------------- GENERAL ----------------
\setlength{\parindent}{0pt}
\pagenumbering{gobble}
\linespread{1.2}

\newcommand{\mychapter}[2]{%
    \phantomsection
    \hypertarget{#2}{}%
    \bookmark[level=0,dest=#2]{#1}%
    \begin{center}
        {\headingfont\Huge #1}
        \par
        \vspace{4pt}
        \hrule
    \end{center}
}

\newcommand{\mysection}[3]{%
    \phantomsection
    \hypertarget{#3}{}%
    \bookmark[level=1,dest=#3]{#1\space #2}%
    \newpage
    \noindent{\headingfont\Large #1\space #2\par}
    \vspace{4pt}
    \hrule
    \vspace{15pt}
}

\newcommand{\mysubsection}[2]{%
    \vspace{5px}
    \noindent{\headingfont\large #1\space #2\par}
    \vspace{4pt}
    \hrule
    \vspace{15pt}
}

\begin{document}

\mychapter{Chapter 2 : Energy, Energy Transfer \& General Energy Analysis}{534y5fdh}

\begin{center}
\renewcommand{\arraystretch}{1.5}
\begin{tabular}{|c|p{14cm}|}
\hline
\textbf{Section} & \textbf{Core Revision Focus} \\
\hline
2.1 & Conservation of energy, energy quality, refrigerator example \\
\hline
2.2 & Forms of energy, total/internal/mechanical energy, KE/PE, flow energy, mass \& energy flow rates \\
\hline
2.3 & Heat transfer, adiabatic process, specific heat transfer, heat-transfer rate, conduction/convection/radiation \\
\hline
2.4 & Work transfer, sign convention, path functions, electrical work, generalized work \\
\hline
2.5 & Mechanical work, shaft/spring/elastic/surface work, raising \& accelerating bodies \\
\hline
2.6 & First law, energy balance, closed/control-volume systems, stationary systems, rate/specific forms, cycles \\
\hline
2.7 & Efficiency, HHV/LHV, combustion/power-plant/mechanical/pump/turbine/motor/generator efficiencies, combined efficiencies \\
\hline
2.8 & Air pollution, smog/ozone, CO, particulate matter, acid rain, greenhouse effect, climate change, $CO_2$ emissions, EVs, renewable energy \\
\hline
\end{tabular}
\end{center}

% ============================================================
\mysection{2.1}{INTRODUCTION}{2.1}

\begin{itemize}
    \item \textbf{Energy conservation:} Total energy is conserved during a process.
    \item \textbf{Energy quality:} High-quality energy (e.g.\ electrical energy) can be completely converted to low-quality thermal energy, but the reverse conversion is limited.
    \item Refrigerator in an insulated room: electrical energy supplied ultimately increases the internal energy and temperature of the room.
\end{itemize}

% ============================================================
\mysection{2.2}{FORMS OF ENERGY}{2.2}

\mysubsection{}{Total Energy}

\begin{equation*}
e=\frac{E}{m}
\end{equation*}

\begin{itemize}
    \item $E$ : Total energy
    \item $e$ : Specific total energy
    \item $m$ : Mass
\end{itemize}

Only changes in energy are required; the reference level for absolute energy may be chosen arbitrarily.

\textbf{Energy classification:}
\begin{itemize}
    \item \textbf{Macroscopic energy:} Associated with the system as a whole relative to an external reference frame.
    \item \textbf{Microscopic energy:} Associated with molecular structure and molecular activity.
    \item Sum of microscopic energies = \textbf{internal energy} $U$.
\end{itemize}

\mysubsection{}{Kinetic Energy}

\begin{equation*}
KE=\frac{mV^2}{2}
\qquad
ke=\frac{V^2}{2}
\end{equation*}

For a rotating solid body:
\begin{equation*}
KE=\frac{1}{2}I\omega^2
\end{equation*}

\mysubsection{}{Potential Energy}

\begin{equation*}
PE=mgz
\qquad
pe=gz
\end{equation*}

$z$ is measured relative to a selected reference level.

\mysubsection{}{Total Energy of a Simple Compressible System}

When magnetic, electric and surface-tension effects are negligible:
\begin{equation*}
E=U+KE+PE
\end{equation*}

\begin{equation*}
E=U+\frac{mV^2}{2}+mgz
\end{equation*}

On a unit-mass basis:
\begin{equation*}
e=u+\frac{V^2}{2}+gz
\end{equation*}

For a stationary system:
\begin{equation*}
\Delta KE=0,
\qquad
\Delta PE=0,
\qquad
\boxed{\Delta E=\Delta U}
\end{equation*}

\mysubsection{}{Internal Energy}

Internal energy includes microscopic kinetic and potential energies.

\begin{itemize}
    \item \textbf{Sensible energy:} Associated with molecular kinetic energy; generally increases with temperature.
    \item \textbf{Latent energy:} Associated with phase of a substance.
    \item \textbf{Chemical energy:} Associated with atomic bonds.
    \item \textbf{Nuclear energy:} Associated with strong binding forces within the nucleus.
\end{itemize}

\mysubsection{}{Mechanical Energy}

Mechanical energy can be completely and directly converted into mechanical work by an ideal mechanical device.

\[
\boxed{\text{Mechanical energy}=
\text{flow}+\text{kinetic}+\text{potential}}
\]

For a flowing fluid:
\begin{equation*}
e_{\mathrm{flow}}=\frac{P}{\rho}
\end{equation*}

\begin{equation*}
e_{\mathrm{mech}}
=
\frac{P}{\rho}
+\frac{V^2}{2}
+gz
\end{equation*}

\begin{equation*}
\dot{E}_{\mathrm{mech}}
=
\dot{m}
\left(
\frac{P}{\rho}
+\frac{V^2}{2}
+gz
\right)
\end{equation*}

For an incompressible fluid:
\begin{equation*}
\Delta e_{\mathrm{mech}}
=
\frac{P_2-P_1}{\rho}
+
\frac{V_2^2-V_1^2}{2}
+
g(z_2-z_1)
\end{equation*}

\begin{equation*}
\Delta\dot{E}_{\mathrm{mech}}
=
\dot{m}
\left[
\frac{P_2-P_1}{\rho}
+
\frac{V_2^2-V_1^2}{2}
+
g(z_2-z_1)
\right]
\end{equation*}

\begin{itemize}
    \item Constant $P,\rho,V,z$ $\Rightarrow$ $\Delta e_{\mathrm{mech}}=0$.
    \item Positive $\Delta e_{\mathrm{mech}}$ $\Rightarrow$ mechanical work supplied to fluid.
    \item Negative $\Delta e_{\mathrm{mech}}$ $\Rightarrow$ mechanical work extracted from fluid.
\end{itemize}

For an ideal turbine:
\begin{equation*}
\dot{W}_{\max}=\dot{m}\Delta e_{\mathrm{mech}}
\end{equation*}

\textbf{Mass flow rate:}
\begin{equation*}
\dot{m}=\rho\dot{V}
\end{equation*}

\begin{equation*}
\dot{V}=A_cV_{\mathrm{avg}}
\end{equation*}

\begin{equation*}
\boxed{\dot{m}=\rho A_cV_{\mathrm{avg}}}
\end{equation*}

\textbf{Energy flow rate:}
\begin{equation*}
\boxed{\dot{E}=\dot{m}e}
\end{equation*}

\begin{equation*}
1\ \mathrm{kJ/s}=1\ \mathrm{kW}
\end{equation*}

% ============================================================
\mysection{2.3}{ENERGY TRANSFER BY HEAT}{2.3}

\mysubsection{}{Heat Transfer}

\textbf{Heat} = energy transferred solely because of a temperature difference.

\begin{itemize}
    \item Heat flows from higher temperature to lower temperature.
    \item No temperature difference $\Rightarrow$ no heat transfer.
    \item Heat is \textbf{energy in transition}; a system does not contain heat as a property.
    \item Heat transfer stops at thermal equilibrium.
\end{itemize}

\mysubsection{}{Adiabatic Process}

\begin{equation*}
\boxed{Q=0}
\end{equation*}

Occurs when:
\begin{itemize}
    \item the system is well insulated, or
    \item system and surroundings are at the same temperature.
\end{itemize}

\textbf{Important:} Adiabatic $\neq$ isothermal. Temperature may change during an adiabatic process because of work or other energy interactions.

\mysubsection{}{Amount of Heat Transfer}

\begin{equation*}
q=\frac{Q}{m}
\end{equation*}

\begin{equation*}
\dot{Q}=\frac{dQ}{dt}
\end{equation*}

For variable heat-transfer rate:
\begin{equation*}
Q=\int_{t_1}^{t_2}\dot{Q}\,dt
\end{equation*}

For constant rate:
\begin{equation*}
Q=\dot{Q}\Delta t,
\qquad
\Delta t=t_2-t_1
\end{equation*}

\begin{equation*}
q=\frac{Q}{m},
\qquad
\dot{Q}=\frac{Q}{\Delta t}
\end{equation*}

\mysubsection{}{Mechanisms of Heat Transfer}

\begin{itemize}
    \item \textbf{Conduction:} Molecular interaction between adjacent particles.
    \item \textbf{Convection:} Energy transfer between a solid surface and moving fluid; includes conduction and fluid motion.
    \item \textbf{Radiation:} Electromagnetic-wave/photon transfer.
\end{itemize}

All require a temperature difference and proceed from higher to lower temperature.

% ============================================================
\mysection{2.4}{ENERGY TRANSFER BY WORK}{2.4}

\mysubsection{}{Work}

Work is energy transfer associated with a force acting through a distance.

\begin{equation*}
w=\frac{W}{m}
\end{equation*}

\begin{equation*}
\dot{W}=\frac{dW}{dt}
\end{equation*}

\textbf{Sign convention:}
\begin{itemize}
    \item Heat \textbf{to} system: positive.
    \item Heat \textbf{from} system: negative.
    \item Work \textbf{by} system: positive.
    \item Work \textbf{on} system: negative.
\end{itemize}

Alternatively, use positive in/out quantities explicitly:
\begin{equation*}
W_{\mathrm{in}}>0,
\qquad
W_{\mathrm{out}}>0
\end{equation*}

A negative calculated value indicates that the actual direction is opposite to the assumed direction.

\mysubsection{}{Heat and Work as Path Functions}

\begin{itemize}
    \item Heat and work exist only while crossing the system boundary.
    \item A system possesses energy, but not heat or work.
    \item Heat/work are associated with a \textbf{process}, not a state.
    \item Their magnitude depends on the path.
    \item Heat and work = \textbf{path functions}.
    \item Thermodynamic properties = \textbf{point functions}.
\end{itemize}

For a property:
\begin{equation*}
\int_1^2 dV=V_2-V_1=\Delta V
\end{equation*}

For work:
\begin{equation*}
\int_1^2\delta W=W_{12}
\neq W_2-W_1
\end{equation*}

$d$ = exact differential of a property; $\delta$ = inexact differential of a path function.

\mysubsection{}{Electrical Work}

\begin{equation*}
W_e=VN
\end{equation*}

\begin{equation*}
\dot{W}_e=VI
\end{equation*}

For variable $V$ and $I$:
\begin{equation*}
W_e=\int_1^2VI\,dt
\end{equation*}

For constant $V,I$:
\begin{equation*}
W_e=VI\Delta t
\end{equation*}

\begin{equation*}
\boxed{
\dot{W}_e=VI=I^2R=\frac{V^2}{R}
}
\end{equation*}

\mysubsection{}{Generalized Work}

\begin{equation*}
\boxed{\delta W=F\,dx}
\end{equation*}

General concept:
\[
\text{Generalized force}\times\text{generalized displacement}
\]

\begin{itemize}
    \item Mechanical: force $\times$ displacement.
    \item Electrical: voltage $\times$ charge.
    \item Magnetic: magnetic field strength $\times$ total magnetic dipole moment.
    \item Electrical polarization: electric field strength $\times$ polarization.
\end{itemize}

% ============================================================
\mysection{2.5}{MECHANICAL FORMS OF WORK}{2.5}

\textbf{Mechanical work:}
\begin{equation*}
W=Fs
\end{equation*}

For variable force:
\begin{equation*}
W=\int_1^2F\,ds
\end{equation*}

Requirements for mechanical work:
\begin{enumerate}
    \item A force must act on the system boundary.
    \item The system boundary must move.
\end{enumerate}

No boundary displacement $\Rightarrow$ no mechanical work.
Expansion into a vacuum $\Rightarrow$ no work because no opposing/driving force exists.

\begin{center}
\renewcommand{\arraystretch}{1.5}
\begin{tabular}{|p{3.0cm}|p{6.0cm}|p{7.0cm}|}
\hline
\textbf{Type} & \textbf{Important relations} & \textbf{Key conditions / quantities} \\
\hline

\textbf{Shaft Work}
&
\begin{equation*}
T=Fr
\end{equation*}
\begin{equation*}
s=2\pi rn
\end{equation*}
\begin{equation*}
W_{\mathrm{sh}}=2\pi nT
\end{equation*}
\begin{equation*}
\dot{W}_{\mathrm{sh}}=2\pi\dot{n}T
\end{equation*}
&
$T$ = torque,\;
$F$ = tangential force,\;
$r$ = moment arm,\;
$n$ = revolutions,\;
$\dot n$ = revolutions/time.

Used for engines, turbines, pumps, compressors.
\\
\hline

\textbf{Spring Work}
&
\begin{equation*}
\delta W_{\mathrm{spring}}=F\,dx
\end{equation*}
\begin{equation*}
F=kx
\end{equation*}
\begin{equation*}
\delta W_{\mathrm{spring}}=kx\,dx
\end{equation*}
\begin{equation*}
W_{\mathrm{spring}}
=
\frac12k(x_2^2-x_1^2)
\end{equation*}
&
Linear elastic spring; $k$ = spring constant.

$x_1,x_2$ measured from the undisturbed position.
\\
\hline

\textbf{Elastic Solid Bar}
&
\begin{equation*}
W_{\mathrm{elastic}}=\int_1^2F\,dx
\end{equation*}
\begin{equation*}
\sigma_n=\frac{F}{A}
\end{equation*}
\begin{equation*}
W_{\mathrm{elastic}}
=
\int_1^2\sigma_nA\,dx
\end{equation*}
&
Within elastic range, a solid bar behaves similarly to a linear spring.

$\sigma_n$ = normal stress.
\\
\hline

\textbf{Liquid Film}
&
\begin{equation*}
W_{\mathrm{surface}}=\int_1^2\sigma_s\,dA
\end{equation*}
\begin{equation*}
dA=2b\,dx
\end{equation*}
\begin{equation*}
F=2b\sigma_s
\end{equation*}
&
$\sigma_s$ = surface tension,\;
$b$ = movable-wire length.

Factor $2$ because a liquid film has two surfaces.
\\
\hline

\textbf{Raising a Body}
&
\begin{equation*}
W=\Delta PE
\end{equation*}
\begin{equation*}
W=mg(z_2-z_1)
\end{equation*}
&
Work transferred to the body equals its increase in gravitational potential energy.
\\
\hline

\textbf{Accelerating a Body}
&
\begin{equation*}
W=\Delta KE
\end{equation*}
\begin{equation*}
W=\frac12m(V_2^2-V_1^2)
\end{equation*}
&
Work required for acceleration equals change in kinetic energy.
\\
\hline

\end{tabular}
\end{center}

% ============================================================
\mysection{2.6}{THE FIRST LAW OF THERMODYNAMICS}{2.6}

\textbf{First law:}
\[
\boxed{\text{Energy in} - \text{Energy out}
=
\text{Change in stored energy}}
\]


Special cases:
\begin{align*}
Q=0
&\Rightarrow
\Delta E= - W_{\mathrm{net, \, out}}
\\
W=0
&\Rightarrow
\Delta E=Q_{\mathrm{net, \, in}}
\end{align*}

\mysubsection{}{General Energy Balance}

\begin{equation*}
\boxed{
E_{\mathrm{in}}-E_{\mathrm{out}}
=
\Delta E_{\mathrm{system}}
}
\end{equation*}

\begin{equation*}
\Delta E_{\mathrm{system}}
=
E_2-E_1
\end{equation*}

For a simple compressible system:
\begin{equation*}
\Delta E
=
\Delta U+\Delta KE+\Delta PE
\end{equation*}

\begin{equation*}
\Delta U=m(u_2-u_1)
\end{equation*}

\begin{equation*}
\Delta KE=
\frac12m(V_2^2-V_1^2)
\end{equation*}

\begin{equation*}
\Delta PE=mg(z_2-z_1)
\end{equation*}

Therefore:
\begin{equation*}
\boxed{
\Delta E
=
m(u_2-u_1)
+
\frac12m(V_2^2-V_1^2)
+
mg(z_2-z_1)
}
\end{equation*}

\mysubsection{}{Stationary Systems}

\begin{equation*}
\Delta KE=0,
\qquad
\Delta PE=0
\end{equation*}

Hence:
\begin{equation*}
\boxed{\Delta E=\Delta U}
\end{equation*}

\mysubsection{}{Energy-Transfer Mechanisms}

Across a system boundary, energy can be transferred by:

\begin{enumerate}
    \item \textbf{Heat}
    \item \textbf{Work}
    \item \textbf{Mass flow}
\end{enumerate}

For a \textbf{closed system}: heat + work only.

For a \textbf{control volume}: heat + work + mass flow.

\textbf{Net heat:}
\begin{equation*}
Q_{\mathrm{net, \, in}}=Q_{\mathrm{in}}-Q_{\mathrm{out}}
\end{equation*}

\textbf{Net work:}
\begin{equation*}
W_{\mathrm{net, \, out}}=W_{\mathrm{out}}-W_{\mathrm{in}}
\end{equation*}

\textbf{Complete energy balance:}
\begin{equation*}
\boxed{
\Delta E_{\mathrm{system}}
=
E_{\mathrm{in}}-E_{\mathrm{out}}
=
(Q_{\mathrm{net, \, in}})
-
(W_{\mathrm{net, \, out}})
+
(E_{\mathrm{mass,in}}-E_{\mathrm{mass,out}})
}
\end{equation*}
\begin{equation*}
\boxed{
\Delta E_{\mathrm{system}}
=
E_{\mathrm{in}}-E_{\mathrm{out}}
=
(Q_{\mathrm{in}}-Q_{\mathrm{out}})
-
(W_{\mathrm{out}}-W_{\mathrm{in}})
+
(E_{\mathrm{mass,in}}-E_{\mathrm{mass,out}})
}
\end{equation*}

Useful eliminations:
\begin{equation*}
\text{Adiabatic: }Q=0
\qquad
\text{No work: }W=0
\qquad
\text{Closed system: }E_{\mathrm{mass}}=0
\end{equation*}

\mysubsection{}{Rate and Specific Forms}

Rate form:
\begin{equation*}
\boxed{
\dot{E}_{\mathrm{in}}-\dot{E}_{\mathrm{out}}
=
\frac{dE_{\mathrm{system}}}{dt}
}
\end{equation*}

For constant rates:
\begin{equation*}
Q=\dot Q\Delta t,
\qquad
W=\dot W\Delta t,
\qquad
\Delta E=\frac{dE}{dt}\Delta t
\end{equation*}

Specific basis:
\begin{equation*}
e=\frac{E}{m}
\end{equation*}

\begin{equation*}
e_{\mathrm{in}}-e_{\mathrm{out}}
=
\Delta e_{\mathrm{system}}
\end{equation*}

Differential form:
\begin{equation*}
\delta e_{\mathrm{in}}-\delta e_{\mathrm{out}}
=
de_{\mathrm{system}}
\end{equation*}

\mysubsection{}{Energy Balance for a Cycle}

For a cycle:
\begin{equation*}
\boxed{\Delta E_{\mathrm{system}}=0}
\end{equation*}

Thus:
\begin{equation*}
E_{\mathrm{in}}=E_{\mathrm{out}}
\end{equation*}

For a closed-system cycle:
\begin{equation*}
\boxed{
W_{\mathrm{net,out}}=Q_{\mathrm{net,in}}
}
\end{equation*}

and in rate form:
\begin{equation*}
\boxed{
\dot W_{\mathrm{net,out}}
=
\dot Q_{\mathrm{net,in}}
}
\end{equation*}

% ============================================================
\mysection{2.7}{ENERGY CONVERSION EFFICIENCIES}{2.7}

\textbf{General definition:}
\begin{equation*}
\boxed{
\eta=
\frac{\text{Desired output}}
{\text{Required input}}
}
\end{equation*}

\begin{equation*}
\eta(\%)=100\eta
\end{equation*}

\textbf{Always define clearly what constitutes the desired output and required input.}

\mysubsection{}{Combustion Equipment}

\begin{equation*}
\boxed{
\eta_{\mathrm{comb,equip}}
=
\frac{Q_{\mathrm{useful}}}
{HV\,m_{\mathrm{fuel}}}
}
\end{equation*}

Applies to furnaces, boilers and heaters.

\textbf{Heating values:}
\begin{itemize}
    \item \textbf{LHV:} Water leaves as vapor; latent heat not recovered.
    \item \textbf{HHV:} Water is condensed; latent heat recovered.
\end{itemize}

\begin{equation*}
\boxed{HHV>LHV}
\end{equation*}

For gasoline:
\begin{equation*}
LHV\approx44{,}000\ \mathrm{kJ/kg}
\end{equation*}

\begin{equation*}
HHV\approx47{,}300\ \mathrm{kJ/kg}
\end{equation*}

\textbf{AFUE} accounts for combustion efficiency plus heat losses, start-up losses and cool-down losses.

\mysubsection{}{Power-Plant Efficiencies}

Thermal efficiency:
\begin{equation*}
\boxed{
\eta_{\mathrm{thermal}}
=
\frac{\dot W_{\mathrm{net,out}}}
{\dot Q_{\mathrm{in}}}
}
\end{equation*}

Generator efficiency:
\begin{equation*}
\boxed{
\eta_{\mathrm{generator}}
=
\frac{\dot W_{\mathrm{elect,out}}}
{\dot W_{\mathrm{shaft,in}}}
}
\end{equation*}

Overall:
\begin{equation*}
\boxed{
\eta_{\mathrm{overall}}
=
\eta_{\mathrm{comb,equip}}
\eta_{\mathrm{thermal}}
\eta_{\mathrm{generator}}
}
\end{equation*}

or
\begin{equation*}
\eta_{\mathrm{overall}}
=
\frac{\dot W_{\mathrm{net,electric}}}
{HHV\,\dot m_{\mathrm{fuel}}}
\end{equation*}

\mysubsection{}{Electricity-to-Light Conversion}

\begin{equation*}
\eta_{\mathrm{light}}
=
\frac{\text{Energy converted to light}}
{\text{Electrical energy consumed}}
\end{equation*}

Lighting efficacy:
\begin{equation*}
\boxed{
\text{Lighting efficacy}
=
\frac{\text{Light output (lumens)}}
{\text{Electrical power input (W)}}
}
\end{equation*}

Unit: $\mathrm{lm/W}$.

\mysubsection{}{Cooking Efficiency}

\begin{equation*}
\boxed{
\eta_{\mathrm{cooking}}
=
\frac{\text{Useful energy transferred to food}}
{\text{Energy consumed by appliance}}
}
\end{equation*}

Convection and microwave ovens are generally more efficient than conventional ovens.

\mysubsection{}{Mechanical Efficiency}

\begin{equation*}
\boxed{
\eta_{\mathrm{mech}}
=
\frac{E_{\mathrm{mech,out}}}
{E_{\mathrm{mech,in}}}
}
\end{equation*}

\begin{equation*}
\eta_{\mathrm{mech}}
=
1-
\frac{E_{\mathrm{mech,loss}}}
{E_{\mathrm{mech,in}}}
\end{equation*}

Losses primarily arise from friction and other irreversibilities.

\mysubsection{}{Pump Efficiency}

\begin{equation*}
\boxed{
\eta_{\mathrm{pump}}
=
\frac{\Delta\dot E_{\mathrm{mech,fluid}}}
{\dot W_{\mathrm{shaft,in}}}
}
\end{equation*}

\begin{equation*}
\eta_{\mathrm{pump}}
=
\frac{\dot W_{\mathrm{pump,u}}}
{\dot W_{\mathrm{pump}}}
\end{equation*}

\mysubsection{}{Turbine Efficiency}

\begin{equation*}
\boxed{
\eta_{\mathrm{turbine}}
=
\frac{\dot W_{\mathrm{shaft,out}}}
{\left|\Delta\dot E_{\mathrm{mech,fluid}}\right|}
}
\end{equation*}

Absolute value is required because fluid mechanical energy decreases through a turbine.

\mysubsection{}{Motor and Generator Efficiency}

Motor:
\begin{equation*}
\boxed{
\eta_{\mathrm{motor}}
=
\frac{\dot W_{\mathrm{shaft,out}}}
{\dot W_{\mathrm{elect,in}}}
}
\end{equation*}

Generator:
\begin{equation*}
\boxed{
\eta_{\mathrm{generator}}
=
\frac{\dot W_{\mathrm{elect,out}}}
{\dot W_{\mathrm{shaft,in}}}
}
\end{equation*}

\[
\text{Motor: Electrical}\rightarrow\text{Mechanical}
\]

\[
\text{Generator: Mechanical}\rightarrow\text{Electrical}
\]

\mysubsection{}{Combined Efficiencies}

Pump--motor:
\begin{equation*}
\boxed{
\eta_{\mathrm{pump-motor}}
=
\eta_{\mathrm{pump}}\eta_{\mathrm{motor}}
}
\end{equation*}

\begin{equation*}
\eta_{\mathrm{pump-motor}}
=
\frac{\dot W_{\mathrm{pump,u}}}
{\dot W_{\mathrm{elect,in}}}
=
\frac{\Delta\dot E_{\mathrm{mech,fluid}}}
{\dot W_{\mathrm{elect,in}}}
\end{equation*}

Turbine--generator:
\begin{equation*}
\boxed{
\eta_{\mathrm{turbine-gen}}
=
\eta_{\mathrm{turbine}}\eta_{\mathrm{generator}}
}
\end{equation*}

\begin{equation*}
\eta_{\mathrm{turbine-gen}}
=
\frac{\dot W_{\mathrm{elect,out}}}
{\dot W_{\mathrm{turbine,e}}}
=
\frac{\dot W_{\mathrm{elect,out}}}
{\left|\Delta\dot E_{\mathrm{mech,fluid}}\right|}
\end{equation*}

% ============================================================
\mysection{2.8}{ENERGY AND ENVIRONMENT}{2.8}

\textbf{Major environmental effects of fossil-fuel combustion:}
\begin{itemize}
    \item Air pollution
    \item Smog
    \item Acid rain
    \item Global warming
    \item Climate change
\end{itemize}

\textbf{Major motor-vehicle pollutants:}
\begin{itemize}
    \item Hydrocarbons (HC)
    \item Nitrogen oxides ($NO_x$)
    \item Carbon monoxide (CO)
    \item Particulate matter (PM)
\end{itemize}

\mysubsection{}{Ozone and Smog}

Ground-level ozone is the primary component of urban smog.

\begin{equation*}
HC+NO_x
\xrightarrow{\text{sunlight}}
O_3+\text{other pollutants}
\end{equation*}

Hot, calm conditions favor smog formation; winds can transport pollutants over large distances.

\mysubsection{}{Carbon Monoxide}

CO is a colorless, odorless, poisonous gas produced mainly by motor vehicles.

\begin{itemize}
    \item Reduces blood's oxygen-carrying ability.
    \item Low concentrations can impair reactions, reflexes and judgment.
    \item High concentrations can be fatal.
    \item Particularly dangerous in enclosed/poorly ventilated spaces.
\end{itemize}

\mysubsection{}{Particulate Matter}

Suspended dust and soot can be emitted by vehicles, industrial facilities and combustion systems.

Effects include irritation of eyes/lungs and transport of harmful substances such as acids and metals.

\mysubsection{}{Acid Rain}

Sulfur combustion:
\begin{equation*}
S+O_2\rightarrow SO_2
\end{equation*}

Atmospheric reactions:
\begin{equation*}
SO_x+H_2O\rightarrow\text{sulfuric acid}
\end{equation*}

\begin{equation*}
NO_x+H_2O\rightarrow\text{nitric acid}
\end{equation*}

These acids reaching the earth through rain/snow produce \textbf{acid rain}.

Reduction methods:
\begin{itemize}
    \item Sulfur-dioxide scrubbers
    \item Low-sulfur coal
    \item Coal desulfurization
\end{itemize}

\mysubsection{}{Greenhouse Effect and Climate Change}

\textbf{Greenhouse effect:} Certain gases absorb outgoing infrared radiation and re-emit part of it toward earth.

Major greenhouse gases:
\begin{itemize}
    \item $CO_2$
    \item Water vapor
    \item Methane
    \item Nitrogen oxides
\end{itemize}

The natural greenhouse effect is necessary for life; excessive greenhouse-gas concentrations increase trapped energy and contribute to global warming.

Major sources:
\begin{itemize}
    \item Power generation
    \item Transportation
    \item Buildings
    \item Manufacturing
    \item Fossil-fuel combustion
\end{itemize}

Photosynthesis:
\begin{equation*}
CO_2+\mathrm{H_2O}
\rightarrow
\text{organic matter}+O_2
\end{equation*}

Deforestation decreases carbon uptake while fossil-fuel use increases $CO_2$ emissions.

\textbf{Consequences:}
\begin{itemize}
    \item Weather-pattern changes
    \item Severe storms, heavy rainfall and flooding
    \item Drought
    \item Polar-ice melting and sea-level rise
    \item Ecosystem disruption
    \item Changes in water availability and disease distribution
    \item Human-health and socioeconomic impacts
\end{itemize}

\textbf{International efforts:}
\begin{itemize}
    \item 1992 United Nations climate agreement
    \item 1997 Kyoto Protocol
    \item 2015 Paris Agreement
\end{itemize}

Paris Agreement: limit warming to less than $2^\circ\mathrm{C}$ above preindustrial levels while pursuing further limitation.

\mysubsection{}{Carbon Dioxide Emissions}

Fossil-fueled power plant:
\begin{equation*}
0.6\text{--}1.0\ \mathrm{kg}\ CO_2/\mathrm{kWh}
\end{equation*}

Gasoline:
\begin{equation*}
\approx2.5\ \mathrm{kg}\ CO_2/\mathrm{L}
\end{equation*}

\begin{equation*}
\approx20\ \mathrm{lbm}\ CO_2/\mathrm{gal}
\end{equation*}

For a vehicle:
\begin{equation*}
\boxed{
CO_2\text{ emission}
=
\text{Fuel consumed}
\times
\text{$CO_2$ emission per unit fuel}
}
\end{equation*}

Therefore:
\begin{equation*}
\boxed{
\text{Higher efficiency}
\rightarrow
\text{Lower fuel consumption}
\rightarrow
\text{Lower emissions}
}
\end{equation*}

\mysubsection{}{Electric Vehicles}

\begin{itemize}
    \item No direct tailpipe emissions during operation.
    \item Not necessarily zero total emissions because electricity generation may involve fossil fuels.
    \item Environmental benefit depends on the electricity source.
\end{itemize}

\begin{equation*}
\text{Electric vehicle}
\neq
\text{necessarily zero total emissions}
\end{equation*}

Emission-free electricity sources listed in the notes:
\begin{itemize}
    \item Hydroelectric
    \item Solar
    \item Wind
    \item Geothermal
\end{itemize}

\mysubsection{}{Renewable Energy and Thermodynamics}

Important renewable sources:
\begin{itemize}
    \item Hydroelectric
    \item Solar
    \item Wind
    \item Geothermal
\end{itemize}

Thermodynamic improvement:
\begin{equation*}
\boxed{
\text{Improved efficiency}
\rightarrow
\text{Less fuel consumption}
\rightarrow
\text{Reduced pollution}
}
\end{equation*}

\textbf{Overall exam chain:}
\[
\boxed{
\text{Energy conversion}
\rightarrow
\text{Efficiency}
\rightarrow
\text{Fuel consumption}
\rightarrow
\text{Environmental impact}
}
\]

\end{document}